A análise dos dados coletados demonstrou que o sistema cumpre sua função básica de monitoramento ambiental. No entanto, é fundamental reiterar que a precisão dos resultados está intrinsecamente ligada às limitações de hardware do sensor utilizado, o DHT11. Com uma margem de erro de $\pm$ 2°C para temperatura e $\pm$ 5\% para umidade, os dados obtidos devem ser interpretados como tendências ambientais e não como medições laboratoriais de alta fidelidade.
\textbf{Considerações sobre o Monitoramento} 

Embora as especificações de operação do DHT11 (0–50°C e 20–95\% de umidade) atendam plenamente às condições climáticas da nossa região, a imprecisão inerente pode ser um fator crítico em aplicações que exijam maior controle de estabilidade, como em estufas ou CPDs.

\textbf{Proposta de Otimização} 

Como diretriz para trabalhos futuros e visando a evolução da robustez do sistema, propõe-se a substituição do módulo atual pelo sensor DHT22 (AM2302). Esta atualização oferece vantagens significativas:
Aumento da Precisão: Redução da margem de erro para $\pm$ 0,5°C em temperatura e $\pm$ 2–5\% em umidade.
Amplitude de Faixa: Capacidade de leitura em escalas negativas e níveis de umidade mais abrangentes.
Confiabilidade: Maior estabilidade nos dados, diminuindo as discrepâncias sazonais e garantindo uma leitura mais próxima do valor real.

\vspace{0.5cm}

Em suma, o protótipo valida a viabilidade do projeto, mas a transição para sensores de maior sensibilidade é o passo natural para garantir que o sistema de monitoramento atinja um patamar profissional de confiabilidade.